\documentclass[10pt,a4paper]{amsart}
\usepackage[margin=19mm]{geometry}
\usepackage{amsmath,amssymb,mathtools}
\usepackage[scr=rsfs,cal=euler]{mathalfa}
\usepackage{xfrac}
\usepackage[unicode,hidelinks,bookmarks=false]{hyperref}
\newcommand{\ie}{i.e.\ }
\newcommand{\cf}{cf.\ }
\newcommand{\resp}{respectively\ }
\newcommand{\Subsection}{\subsection}
\providecommand{\nobreakdash}{\nobreak\hbox{-}}
\setlength{\emergencystretch}{2em}
\multlinegap0pt
\usepackage{bm}
\usepackage{bbm}
\usepackage{dsfont}
\usepackage{xspace}
\usepackage{cancel}
\usepackage{mathtools}
\usepackage{amsthm}
\makeatletter
\@ifundefined{theorem}{\newtheorem{theorem}{Theorem}}{}
\makeatother
\makeatletter
\expandafter\ifx\csname remark\endcsname\relax
\newenvironment{remark}[1][Remark]{\begin{trivlist}
\item[\hskip \labelsep {\bfseries #1}]}{\end{trivlist}}
\fi
\makeatother
\title[Discrete-Time Two-Strain Epidemic Dynamics on Complex Networ]{Discrete-Time Two-Strain Epidemic Dynamics on Complex Networks}
\author{Frank Namugera}
\date{Source: arXiv:2508.08294v2 (CC BY 4.0)}
\begin{document}
\maketitle

\noindent\emph{Source and licence.} Discrete-Time Two-Strain Epidemic Dynamics on Complex Networks. Version arXiv:2508.08294v2 (CC BY 4.0), distributed under the Creative Commons Attribution 4.0 International licence, as recorded in the arXiv licence metadata for this exact version and shown on the abstract page https://arxiv.org/abs/2508.08294v2. Full rights notice: licenses/arxiv-2508.08294-CC-BY-4.0.txt.\par\smallskip

\noindent\textit{Editorial scope.} Counted unit: the Theorem whose source label is \texttt{thm:threshold} in the article (source0.tex lines 316--329, source label \texttt{thm:threshold}), delivered together with the article's own complete original proof of it (source0.tex lines 331--376, one \texttt{proof} environment of 46 lines). No mathematics is written, repaired or re-derived: statement and proof are the article's own text line for line. Required context retained uncounted: thm:pefro (lines 812--820). Every definition, display and label of those lines is retained unchanged. Omitted from this package and not redistributed: 1--315, 330--330, 377--811, 821--830. Nothing inside a retained excerpt is omitted or altered: every retained body is delivered exactly as the article's lines are written, so no omission receipt of any kind accompanies this package. Citations: the retained excerpts carry no citation command at all, so no bibliography entry of the article is transcribed and no reference is redistributed. Reference adaptation: none was needed and none was made; every reference inside the delivered unit resolves inside it, so no unresolved reference or citation is produced. Printed slips kept literal: no printed word, symbol, hypothesis or number of the counted statement or of its proof was altered. Layout and class: the article's own class is \texttt{scrartcl} with packages including amsmath, amssymb, amsthm, stmaryrd, bbm, hyperref, mathtools, natbib, enumitem, subcaption, float, authblk, tikz, tikz-3dplot; the wrapper is the lane's pinned \texttt{amsart} editorial wrapper at 10pt on A4, which restates the theorem-like environments the retained text needs and adds the article's own macro definitions that the retained text needs (none), with extra wrapper packages (bm, bbm, dsfont, xspace, cancel, mathtools). The delivered numbering is the unit's own. Assets: no retained span contains a figure environment, a graphics-inclusion command, a PDF-page inclusion, a table or a TikZ picture; no image file is copied into this package, the package declares no asset dependency and no image file is redistributed with it. Licence: version arXiv:2508.08294v2 of this article is distributed under the Creative Commons Attribution 4.0 International licence, as recorded in the arXiv licence metadata for this exact version and shown on the abstract page https://arxiv.org/abs/2508.08294v2. A full rights notice is delivered at licenses/arxiv-2508.08294-CC-BY-4.0.txt. Characters: every retained excerpt is pure ASCII, so the delivered engine, XeLaTeX with its default Latin Modern text font, reports no missing character.
\begin{theorem}[Perron-Frobenius]
Let $A \in \mathbb{R}^{n \times n}$ be a nonnegative and irreducible matrix. Then:  
\begin{enumerate}
    \item $A$ has a simple eigenvalue $\lambda_1 > 0$, such that $\lambda_1 = \max\{|\lambda| : \lambda \in \rho(A)\}$, where $\rho(A)$ denotes the spectrum of $A$.
    % \item There exists a strictly positive eigenvector $v \in \mathbb{R}^n_{>0}$ such that $Av = \rho(A)v$.
    \item Every other eigenvalue $\lambda_i$ of $A$ satisfies $|\lambda| \leq \lambda_1$.
\end{enumerate}
\label{thm:pefro}
\end{theorem}

\begin{theorem}[Epidemic Threshold and Local Stability]
\label{thm:threshold}
% Let 
% $
% W_K := \bigl(R \circ D^{-\alpha} \circ \mathbf{K}\bigr) \in \mathbb{R}^{n \times n},
% $
% be the infection matrix associated with strain $K \in \{A,B\}$, where $W_K$ has nonnegative entries and is irreducible.  
Let $\lambda_1(W_K)$ be the spectral radius of $W_{K}$, and consider the linearization of the system at the disease-free equilibrium $\mathbf{p}^* = \mathbf{0} \in \mathbb{R}^{3n}$.  
Then $\mathbf{p}^*$ is locally asymptotically stable if 
\[
\frac{\sigma_K}{\delta} < \frac{1}{\lambda_{1}(W_K)}
\]
and is otherwise unstable for $K \in \{A,B\}$.
\end{theorem}

\begin{proof}

We study the stability of the disease-free configuration 
\(\mathbf{p}^* = \mathbf{0}\) 
under the discrete-time evolution
\[
\mathbf{p}_{t+1} = g(\mathbf{p}_t).
\]
The linearized dynamics near $\mathbf{p}^*$ are governed by the Jacobian $(J)$. 
% \[
% J := \nabla g(0).
% \]

By construction, $J$ is block upper-triangular, with diagonal blocks
\[
J_A = (1-\delta) \mathrm{I}_n + \sigma_{A} W_A, 
\qquad
J_B = (1-\delta) \mathrm{I}_n + \sigma_{B} W_B,
\qquad
J_{AB} = (1-\mu)\mathrm{I}_n,
\]
where $J_A$ and $J_B$ act on the infection probabilities of strains $A$ and $B$, and $J_{AB}$ corresponds to the co-infection component. Since the spectrum of a block upper-triangular matrix is the union of the spectra of its diagonal blocks, then
\[
\operatorname{spec}(J) = \operatorname{spec}(J_A) \cup \operatorname{spec}(J_B) \cup \operatorname{spec}(J_{AB}) .
\]

For discrete-time dynamics, linear stability requires that all eigenvalues of $J$ lie strictly inside the unit disk, i.e.
\[
\rho(J) < 1,
\]
The scalar eigenvalue $1-\mu$ satisfies $|1-\mu| < 1$ whenever $\mu>0$, so it does not influence the stability threshold. Hence the relevant contributions come from the infection blocks $J_A$ and $J_B$.

Since $W_K$ is nonnegative and irreducible, the Perron--Frobenius theorem (\ref{thm:pefro}) implies that its spectral radius is a real, simple eigenvalue $\lambda_{1}(W_K)$. Thus the leading eigenvalue of each block is shifted by the survival term $(1-\delta)$:
\[
\rho(J_K) = (1-\delta) + \lambda_{1}(\sigma_{K} W_K), \qquad K \in \{A,B\}.
\]

Therefore, the disease-free equilibrium is locally asymptotically stable if 
\[
\sigma_{A} \lambda_{1}(W_A) < \delta
\quad \text{and} \quad
\sigma_{B} \lambda_{1}(W_B) < \delta.
\]

If either condition fails, the spectral radius of $J$ exceeds one, leading to exponential growth of small perturbations, and the disease-free state becomes linearly unstable.
\end{proof}

\end{document}
